\section{Gap analysis}
\label{sec:gap}

\begin{keyfinding}{The gap, in one sentence}
Bhoonidhi stores the footprint polygon of every scene it holds, and the only
question you are permitted to ask of that geometry is \emph{``does it touch
mine?''}
\end{keyfinding}

\noindent Everything in this section is either evidence for that sentence or
an argument about why it matters more than it sounds like it does.

\subsection{What is missing --- the capability audit}

Method 1 (Section~\ref{sec:methodology}) produces
Table~\ref{tab:audit}, written to \code{docs/evidence/capability-audit.json}.

\begin{table}[htbp]
\centering\small
\caption{Bhoonidhi's spatial and analytical surface, audited against its own
published specification.}
\label{tab:audit}
\begin{tabularx}{\textwidth}{@{}>{\raggedright\arraybackslash}p{0.44\textwidth}>{\raggedright\arraybackslash}X@{}}
\toprule
\textbf{Question} & \textbf{Answer} \\
\midrule
Spatial parameters on \ep{/data/search}
  & \code{bbox} and \code{intersects} --- and nothing else \\
What they compute
  & the same single predicate: \emph{does this item intersect my AOI?} \\
Other spatial operators documented
  & \absent\ --- no \code{crosses}, \code{touches}, \code{overlaps},
    \code{disjoint}, \code{s\_within}, \code{s\_contains} \\
CQL2 operations in the documented examples
  & \code{eq}, \code{and} --- attribute equality only \\
Analysis vocabulary anywhere in the API
  & \absent\ --- no union, aggregation, statistics, coverage, density \\
Endpoints documented
  & 5: authenticate, logout, list, search, download \\
Endpoints that compute anything
  & \textcolor{bad}{\textbf{0}} \\
\bottomrule
\end{tabularx}
\end{table}

So the catalogue holds the geometry and will not compute with it. The
questions that actually decide an acquisition order are all exactly one step
beyond what it offers (Table~\ref{tab:questions}).

\begin{table}[htbp]
\centering\small
\caption{Ordinary spatial questions, the operations they require, and their
availability on Bhoonidhi today. Every input these need is geometry the
catalogue already stores.}
\label{tab:questions}
\begin{tabularx}{\textwidth}{@{}>{\raggedright\arraybackslash}X l c@{}}
\toprule
\textbf{Question an analyst asks} & \textbf{Operations required} & \textbf{Today} \\
\midrule
How much of my AOI do these scenes really cover, and where are the holes?
  & union, difference, area & \no \\
Which few of the 30 scenes returned do I actually need?
  & set cover over geometry & \no \\
Where do optical and radar co-cover, for fusion?
  & intersection of two unions & \no \\
How uneven is revisit across my AOI?
  & gridding, per-cell counts & \no \\
What terrain is under each candidate scene?
  & raster--vector overlay & \no \\
\bottomrule
\end{tabularx}
\end{table}

\subsection{This gap is not about authentication}
\label{sec:notauth}

Stated explicitly and early, because it is the most tempting \emph{wrong}
version of this argument.

Requiring a token is normal, sensible practice. Copernicus, USGS
EarthExplorer and NASA Earthdata all gate downloads. This project makes no
complaint about it whatsoever.

\begin{keyfinding}{The gap stands for a fully authenticated user}
Hand someone a valid JWT and the complete API documentation, and every
limitation in Table~\ref{tab:audit} is still in force. They can learn
\emph{which} scenes touch their area, and nothing further. This is a
\textbf{capability} gap, not an \textbf{access} gap.
\end{keyfinding}

\subsection{Why this is a large gap, not a small one}

Four independent arguments. Each is measured rather than asserted.

\subsubsection{Argument 1 --- an entire standardised tier is absent}

Server-side geospatial computation is neither exotic nor speculative. It has
been an OGC standard since 2007 (Table~\ref{tab:standards}).

\begin{table}[htbp]
\centering\small
\caption{Standards for server-side geospatial processing, and Bhoonidhi's
implementation status.}
\label{tab:standards}
\begin{tabular}{@{}llll@{}}
\toprule
\textbf{Standard} & \textbf{Document} & \textbf{Year} & \textbf{Bhoonidhi} \\
\midrule
OGC Web Processing Service 1.0.0 & 05-007r7 & 2007 & \no \\
OGC WPS 2.0.0 & 14-065 & 2015 & \no \\
\textbf{OGC API --- Processes, Part 1: Core} & \textbf{18-062r2} & \textbf{2021} & \no \\
openEO API (EO process graphs) & community & 2020 & \no \\
\bottomrule
\end{tabular}
\end{table}

Method 2 probes ISRO's public stack for any of it:

\begin{lstlisting}
Bhuvan WPS 1.0.0 GetCapabilities       HTTP 404   absent
Bhuvan OGC API - Processes landing     HTTP 404   absent
Bhoonidhi /processes                   HTTP 404   absent
  MATCHED  no server-side computation anywhere in the public stack
\end{lstlisting}

The missing thing is therefore not a query parameter someone forgot. It is
the whole processing tier of the OGC stack --- the tier that turns a catalogue
into something you can compute with --- absent from India's national remote
sensing portal nearly two decades after it was standardised.

\subsubsection{Argument 2 --- every comparable portal has it}

The strongest objection to Argument 1 is that the standard is a paper
standard nobody deploys. Method 2 tests this against the field, live
(Table~\ref{tab:peers}).

\begin{table}[htbp]
\centering\small
\caption{Processing capability at comparable Earth observation portals,
probed live and unauthenticated. The lower two rows are ISRO's.}
\label{tab:peers}
\begin{tabularx}{\textwidth}{@{}>{\raggedright\arraybackslash}p{0.33\textwidth}>{\raggedright\arraybackslash}X r@{}}
\toprule
\textbf{Portal} & \textbf{Endpoint probed} & \textbf{Result} \\
\midrule
\textbf{Copernicus Data Space} (openEO)
  & \ep{/openeo/1.2/processes}
  & \textcolor{good}{\textbf{200 --- 146 processes}} \\
MS Planetary Computer
  & STAC \ep{/api/stac/v1/}
  & \textcolor{good}{200 --- 15 classes} \\
\midrule
\textbf{Bhoonidhi}
  & \ep{/processes}
  & \textcolor{bad}{\textbf{404}} \\
\textbf{Bhuvan}
  & \ep{/bhuvan/wps}
  & \textcolor{bad}{\textbf{404}} \\
\bottomrule
\end{tabularx}
\end{table}

Documented but deliberately not probed, each requiring registration:
\textbf{NASA Earthdata Harmony}~\cite{harmony} (server-side subsetting,
reprojection, regridding and aggregation as a service), \textbf{USGS ESPA /
M2M} (on-demand higher-level Landsat products), \textbf{Google Earth
Engine}~\cite{gee} (computation \emph{is} the access model), and
\textbf{Digital Earth Africa} (Open Data Cube with OGC API --- Processes).

\begin{keyfinding}{The contrast}
The Copernicus Data Space publishes \textbf{146} server-side processes at a
public, unauthenticated endpoint that anyone can enumerate right now.
Bhoonidhi publishes \textbf{0}. Bhuvan's WPS endpoint returns \textbf{404}.

\medskip
That is not a scoping decision that everyone else also made. It is the single
capability that separates Bhoonidhi from every portal it would be compared
against.
\end{keyfinding}

\subsubsection{Argument 3 --- the cost is measurable, and it is 92\,\%}
\label{sec:cost}

This is the decisive argument.

Because the catalogue will not compute with footprint geometry, a user has no
way to learn which of the returned scenes they actually need. So the rational
move --- and the observed one --- is to order everything the search returned.
Method 3 measures how much of that is waste (Table~\ref{tab:redundancy}).

\begin{table}[htbp]
\centering\small
\caption{Measured redundancy. \emph{found} is what a search hands back;
\emph{needed} is what covers the same ground to the same percentage, computed
by this service's \code{minimal-scene-set}. Scene counts are measured; GB
figures are published nominal product sizes multiplied by those counts.}
\label{tab:redundancy}
\footnotesize
\begin{tabular}{@{}llrrrrrr@{}}
\toprule
\textbf{AOI} & \textbf{Collection} & \textbf{km\textsuperscript{2}} &
\textbf{found} & \textbf{needed} & \textbf{redundant} &
\textbf{GB today} & \textbf{GB needed} \\
\midrule
taluk             & sentinel-2-l2a &    479 & 30 & 1 & 97\,\% & 30.0 & 1.0 \\
taluk             & sentinel-1-grd &    479 &  4 & 1 & 75\,\% &  3.7 & 0.9 \\
district          & sentinel-2-l2a &  4\,191 & 30 & 2 & 93\,\% & 30.0 & 2.0 \\
district          & sentinel-1-grd &  4\,191 &  8 & 2 & 75\,\% &  7.4 & 1.9 \\
two districts     & sentinel-2-l2a & 11\,855 & 60 & 2 & 97\,\% & 60.0 & 2.0 \\
two districts     & sentinel-1-grd & 11\,855 &  8 & 2 & 75\,\% &  7.4 & 1.9 \\
regional transect & sentinel-2-l2a & 59\,816 & 60 & 4 & 93\,\% & 60.0 & 4.0 \\
regional transect & sentinel-1-grd & 59\,816 &  8 & 2 & 75\,\% &  7.4 & 1.9 \\
\midrule
\textbf{Total} & & & \textbf{208} & \textbf{16} & \textbf{92\,\%} &
\textbf{206.0} & \textbf{15.5} \\
\bottomrule
\end{tabular}
\end{table}

\begin{keyfinding}{The measured cost of the missing tier}
Across 8 queries over 4 areas of interest the catalogue hands back
\textbf{208 scenes}. \textbf{16} of them cover the same ground to the same
percentage. \textbf{192 are redundant --- 92\,\% of everything returned.}

\medskip
At published product sizes that is \textbf{206\,GB ordered where 16\,GB was
required}: 190\,GB of avoidable transfer, for four areas of interest.
Projected onto Bhoonidhi's own ResourceSat-2 LISS-III volumes (250\,MB per
scene), 47\,GB.

\medskip
\textbf{And nothing in Bhoonidhi's API lets a user discover the 92\,\%
figure.}
\end{keyfinding}

As set out in Section~\ref{sec:methodology}, this is a \textbf{floor}: several
rows are capped by the depth of the footprint store rather than by the area of
interest, and redundancy rises with archive depth. A multi-decade archive
returns far more intersecting scenes for the same AOI while the number needed
to cover it stays roughly constant.

\subsubsection{Argument 4 --- the workaround is expert-only, so nobody does it}

The honest counter-argument is \emph{``you can just do this yourself.''} True
--- and that is precisely the gap. The complete procedure a user must follow
today to learn the 92\,\% figure for their own AOI:

\begin{enumerate}
  \item authenticate, \code{POST /data/search}, page through every result;
  \item extract each item's footprint geometry from the STAC response;
  \item load 30+ polygons into a desktop GIS or a \code{shapely} script;
  \item union them, then difference the union against the AOI to find holes;
  \item reproject to an equal-area CRS --- or every area number is wrong;
  \item implement a \textbf{set-cover} optimisation over the footprints;
  \item map the chosen geometries back to orderable item identifiers.
\end{enumerate}

Step 6 is an NP-hard combinatorial problem~\cite{setcover}. Steps 4--5 are
where most people silently get it wrong, because areas computed in degrees are
meaningless. This is a GIS-literate programmer's afternoon, repeated for every
area of interest, by every user, indefinitely --- and it is the \emph{same
seven steps every time}, over geometry the server already holds in memory at
the moment it answers the search.

So in practice nobody does it. They download all 30 scenes. This service
performs all seven steps in one HTTP request (Section~\ref{sec:results}).

\subsection{Counter-arguments, answered}

Taken seriously rather than ignored, because a gap claim that cannot survive
its own objections is not convincing.

\begin{description}[leftmargin=0pt,style=unboxed,itemsep=0.6ex]
  \item[``Requiring a token is the real barrier.''] No ---
    Section~\ref{sec:notauth}. Every limitation here applies to a fully
    authenticated user, and gating downloads is normal practice.
  \item[``The data isn't secret, so there's no gap.''] Agreed, and the gap is
    not about secrecy. It is that there is no way to \emph{compute} with the
    geometry. Section~\ref{sec:cost} measures what that costs.
  \item[``Server-side processing is a paper standard.''] Measured against the
    field in Table~\ref{tab:peers}: Copernicus publishes 146 processes
    publicly, today, at an endpoint this project enumerates live.
  \item[``Users can script it themselves.''] They can, in seven steps
    including an NP-hard optimisation and a reprojection that is easy to get
    silently wrong. The purpose of a service interface is that the server does
    once, correctly, what every client would otherwise do repeatedly and
    inconsistently.
  \item[``STAC search \emph{is} the spatial API.''] STAC search is a
    \emph{filter}: it selects items. It does not perform geometry operations
    and the specification does not claim to. That is exactly the division of
    labour OGC API --- Processes exists to fill.
  \item[``This is a niche optimisation.''] 92\,\% of returned scenes are
    redundant, and the figure rises with archive depth. Order-of-magnitude
    reductions in transfer volume are not niche, particularly for users on
    Indian academic and state-government bandwidth, who are the portal's
    primary audience.
\end{description}

\subsection{Scope of the claim}

\textbf{Claimed:} Bhoonidhi has no server-side spatial computation; that
capability is standardised, deployed by peers, and its absence has a
measurable cost; and a conformant OGC API --- Processes service over footprint
geometry closes it.

\textbf{Not claimed:} that Bhoonidhi is badly built (Section~\ref{sec:gap} is
explicit that discovery is solved), that authentication is wrong, that this
project's footprints \emph{are} Bhoonidhi's (they are substituted and
flagged --- Section~\ref{sec:limitations}), or that this is an official NRSC
product.
